\documentclass{article}
\usepackage{math_template}

\author{Jonathan Kasongo}
\title{Serbia Team Selection Test Round 1 2005 --- P3}

\begin{document}
\maketitle

\begin{problem}{Serbia Team Selection Test Round 1 2005 --- P3}
Find all polynomials $P\left(x\right)$ that satisfy
$P\left(x^2+1\right) = P\left(x\right)^2 + 1$ for all $x$.
\end{problem}

\begin{solution}{Solution}
\textit{Lemma: Any polynomial $P(x) \neq x$ has finitely many fixed
points.}\\

\textit{Proof:}
Suppose not, then $P(x) - x$ has infinitely many solutions,
but since the degree of the polynomial $P(x) - x$ is finite this is a
contradiction. $\square$ \\

Clearly $P(x) = x$ is a solution, now let's find polynomials $P(x) \neq x$
that work also. \\

\textit{Claim: If $t$ is a fixed point of the polynomial $P$, then so
is $t^2 + 1$.}\\

\textit{Proof:} By definition of a fixed point and using the given equation
gives $P\left(t^2 + 1\right) = P\left(t\right)^2 + 1 = t^2 + 1$ and so
$t^2 + 1$ is a fixed point. $\square$ \\

Now suppose that some complex number $z$ is a fixed point of $P$. We must
have $|z| = |z^2 + 1|$ in order to avoid having infinitely many fixed
points. If we write $z = a + bi$ for some real $a,b$ then that last
equation is equivalent to $a^2 + b^2 = (a^2 - b^2 + 1)^2 + (2ab)^2$.
If we think of this equation as a polynomial of, say $a$, the polynomial
is of degree 4 and there are at most 4 solutions. The 4 solutions are
$z = \frac{1}{2} \pm \frac{\sqrt{3}i}{2},
-\frac{1}{2} \pm \frac{\sqrt{3}i}{2}$ and these are easily verifiable
since they correspond to the roots of $x^2 - x + 1 = 0$ and
$x^2 + x + 1 = 0$. Since these are the only fixed points of $P$ we can
write that
$
P(x) - x = (x^2 - x + 1)(x^2 + x + 1)
$
or
$
P(x) - x = x^2 - x + 1
$
or
$
P(x) - x = x^2 + x + 1
$
since we know that roots come in complex conjugate pairs. Solve for
$P(x)$ in each of these equations to show that the candidates for $P(x)$
are $x^4 + x^2 - x + 1$, $(x+1)^2$ and $x^2 + 1$. Then finally one can
check that the only solutions are $P(x) = x \ \text{or}\
  x^2 + 1$.\\

\textit{Remark: Note that this solution is actually only partially
correct. Instead of the condition $|z| = |z^2 + 1|$ the complete condition
is actually that
the sequence $z, \ z^2 + 1, \ (z^2 + 1)^2 + 1, ...$ must have some finite
period, and this means that there are finitely many fixed points as we
initially
desired. Once you work this logic out the actual answer is
$P(x) = x \ \text{or} \ x^2 + 1$ or $x^2 + 1$ composed with itself any
positive integer number of times.}
\end{solution}

\end{document}
